\subsection{Endpoint Principle for the Original Yield}
\label{app:original-yield}

We prove the curvature bound
\begin{equation}
 \begin{split}
 &\frac{-c''(v)[vJ'(v)-J(v)]}{J''(v)}-[c(v)-vc'(v)]<\frac5{32},\\
 &\hspace{12em}v>0.
 \end{split}
 \label{eq:original-shape}
\end{equation}
It implies that \(Y_P(v)=[r c(v)+A_0+A_1J(v)]/v\), with
\(A_0,A_1>0\), has no interior maximum if \(A_0/r\ge5/32\):
its derivative has the sign of
\[
 A_1-\frac{A_0+r[c(v)-vc'(v)]}{vJ'(v)-J(v)},
\]
and the fraction strictly decreases by \eqref{eq:original-shape}.
For the Parseval yield, \(A_0=pd_2/2\) and
\(A_1=p(d_1-d_2/2)F(\ell)\); Appendix~\ref{app:canonical-prerequisites}
verifies \(A_0/r>3/16>5/32\) for \(\ell\ge7/2\).

\paragraph{Entropy curvature.}
Write \(r_v=\tanh v\), \(\delta_v=1-r_v^2\), and
\(D=vJ'-J>0\). Chebyshev's integral inequality gives
\[
 \frac{h(r_v)}{\delta_v}
 =\int_0^1\frac{ds}{(1+s)(1-r_v^2s^2)}
 \le L\int_0^1\frac{ds}{1-r_v^2s^2}=\frac{Lv}{r_v}.
\]
For \(0<v\le1\), this is at most \(L\coth1<1\), since
\(e^2>7\) and \(L<3/4\). Thus
\({\mathcal H}''=(h(r_v)-\delta_v)/\sqrt{\delta_v}<0\).
Consequently \({\mathcal H}^{-2}\) is increasing and convex, and
\(J'''=(\Lambda{\mathcal H}^{-2})''>0\). Hence
\begin{equation}
 D=\int_0^v tJ''(t)\,dt<\frac{v^2}{2}J''(v)
 \qquad(0<v\le1).
 \label{eq:original-profile-curvature}
\end{equation}
We extend this to
\begin{equation}
 \frac{D}{J''}<\frac v2\qquad(v>0).
 \label{eq:profile-ratio-linear}
\end{equation}
Only \(v\ge1\) remains. Put
\(q=e^{-2v}\), \(T=\log(1+q)/q\), and \(K=2v+(1+q)T\).
Then
\begin{gather*}
 J=\frac{1-q}{qK},\qquad K'=2T,\\
 qD=1-\frac{1-q+2T}{K}+\frac{(1-q^2)T^2}{K^2},\\
 0<T'=2[T-1/(1+q)]<q.
\end{gather*}
The last bound is the trapezoid estimate for \(1/(1+q)\).
Using \(q<1/7\), \(1-q/2\le T<1\), \(K>2\), and
\[
 qK<e^{-2v}(2v+3/2)\le(7/2)e^{-2}<1/2,
\]
and dropping the positive derivative of \((1-q^2)T^2\), we obtain
\begin{align*}
 (qD)'\ge\frac1{K^2}\biggl[&2T(1-q+2T)-K(2q+2T')\\
 &-\frac{4(1-q^2)T^3}{K}\biggr]>0:
\end{align*}
the three terms are respectively greater than \(4\), less than
\(2\), and less than \(2\).
Since \((qD)'=q(vJ''-2D)\), this proves
\eqref{eq:profile-ratio-linear}.

\paragraph{Angular curvature.}
Put \(w=\arcsin(\tanh v)/\sinh v\), \(\tau_v=\sech v\),
and \(k=-c''>0\). The angle bound
\eqref{eq:middle-angle-brackets} gives
\(w\ge3\tau_v/(2+\tau_v)>\tau_v^2\).
Using \(w'=(\tau_v^2-w)/r_v<0\), differentiation yields
\[
 vw''-w'\ge
 \frac{\tau_v(1-\tau_v)}{r_v^2(2+\tau_v)}
 \bigl[(3+\tau_v)r_v+v\{4-(1+\tau_v)^3\}\bigr]>0.
\]
Indeed, \(r_v\ge v\tau_v\) bounds the bracket below by
\(v(3-2\tau_v^2-\tau_v^3)>0\).
Thus \(f=1-w\) is increasing and concave in \(v^2\), with
\(f(0)=0\). For \(0<t\le v\),
\(f(t)\ge(t/v)^2f(v)\) and \(f'(t)\ge(t/v)f'(v)\).
Since \(k=2ff'\),
\begin{equation}
 c(v)-vc'(v)=\int_0^v tk(t)\,dt\ge\frac{v^2}{5}k(v).
 \label{eq:angle-curvature-scaling}
\end{equation}
Maximizing the quadratic in
\(k=2(1-w)(w-\delta_v)/r_v\) gives
\(k\le r_v^3/2<1/2\). Therefore
\[
 \frac{kD}{J''}-(c-vc')
 <k(v)\left(\frac v2-\frac{v^2}5\right)\le\frac5{32},
\]
which proves \eqref{eq:original-shape}.
